% ===================================================================
% preamble.tex — Praeambel des vertieften Filters Dual Champions
% ===================================================================
\usepackage[T1]{fontenc}
\usepackage[utf8]{inputenc}
\usepackage{lmodern}
\usepackage[ngerman]{babel}
\usepackage{csquotes}
% patch ohne "footnote": microtype richtet das Fussnotenzeichen im Rand aus,
% kann seinen Patch aber nicht anlegen, wenn footmisc (para) den
% Fussnotenapparat ersetzt - ohne die Einschraenkung meldet jeder Lauf
% "Unable to apply patch `footnote'". Die uebrigen Patches bleiben aktiv.
\usepackage[patch={item,toc,eqnum}]{microtype}
\usepackage{xcolor}
\usepackage{amsmath}
% amssymb liefert \blacktriangle und \blacktriangledown - die Symbole, mit
% denen die Bildunterschrift von fig:kursverlauf die Markenformen benennt.
% Damit braucht die Abbildung keine Legende, die sonst in die Kurve ragt.
\usepackage{amssymb}
\usepackage{booktabs}
\usepackage{longtable}
\usepackage{array}
\usepackage{tabularx}
\usepackage{siunitx}
\usepackage{xfp}
\usepackage{graphicx}
% pgfplots bindet tikz zwar implizit ein, ein eigener Ladeaufruf macht die
% Abhaengigkeit robust gegen eine spaetere Entfernung von pgfplots.
\usepackage{tikz}
% Haelt Gleitobjekte in ihrem Abschnitt: Ohne \FloatBarrier wandert eine
% Tabelle, die den gesamten Rumpf einer Unterunterabschnitts bildet, auf die
% naechste Seite - die Ueberschrift bleibt dann ohne Inhalt zurueck.
\usepackage{placeins}
% Haelt eine Ueberschrift mit dem, was ihr folgt, zusammen. LaTeX verhindert
% von sich aus nur den Umbruch unmittelbar NACH einer Ueberschrift
% (\@afterheading setzt clubpenalty hoch); eine Ueberschrift plus eine
% einzelne Zeile am Seitenfuss bleibt erlaubt - und genau das ist zweimal
% aufgetreten. \needspace{n\baselineskip} vor der Ueberschrift schiebt sie
% auf die naechste Seite, wenn dort nicht mehr n Zeilen Platz haben.
\usepackage{needspace}
\usepackage{pgfplots}
\usepackage{pgfplotstable}
\usepackage[hidelinks]{hyperref}
\usepackage{url}

% --- Satzspiegel: Seitenraender schmaler --------------------------------
% DIV=12 ergibt bei A4 links und rechts je 74.69pt (26.25mm) Rand. Beide
% Raender werden um 10% verkuerzt (auf 67.22pt / 23.6mm); die Textbreite
% waechst dadurch von 448.13pt auf 463.07pt. Die Texthoehe (636.6pt, ein
% Vielfaches des Zeilenabstands) bleibt unveraendert, damit die
% Grundlinien und der Fussnotenapparat unangetastet bleiben.
% \areaset ist der KOMA-eigene Weg; das geometry-Paket wuerde die
% typearea-Berechnung uebersteuern.
\areaset[current]{463.07pt}{636.6pt}

% --- Schusterjungen und Hurenkinder verbieten -----------------------
% LaTeX-Vorgabe ist je 150 und damit so niedrig, dass eine einzelne Zeile
% eines Absatzes allein am Seitenfuss (Schusterjunge) oder am Seitenkopf
% (Hurenkind) stehen bleiben darf. Im deutschen Satz gelten beide als
% Fehler. 10000 verbietet sie vollstaendig; die entstehende Differenz
% faengt die Seitenhoehe auf, da KOMA hier nicht auf gleiche Satzhoehe
% zwingt. \displaywidowpenalty deckt denselben Fall vor einer abgesetzten
% Formel ab - davon gibt es in Teil 2 mehrere.
\clubpenalty=10000
\widowpenalty=10000
\displaywidowpenalty=10000

% --- Fussnotenapparat flach halten ----------------------------------
% Die Zitierweise setzt jede Quelle in eine Fussnote; untereinander gesetzt
% belegte der Apparat auf belegdichten Seiten bis zu einem Drittel des
% Satzspiegels. footmisc mit "para" setzt alle Fussnoten einer Seite in einen
% durchlaufenden Absatz ueber die volle Breite - dieselbe Ersparnis wie ein
% zweispaltiger Apparat, ohne dessen Paketkonflikt: das zweispaltige
% dblfnote (yafoot) laeuft mit pgfplots und hyperref zusammen in eine
% Endlosschleife der Ausgaberoutine (Seiten werden ohne Ende ausgegeben,
% "TeX capacity exceeded, input stack size"); reproduziert mit article und
% scrartcl, unabhaengig von der Ladereihenfolge.
\usepackage[para]{footmisc}

% --- Kopfzeile: Abschnitt links, Dokumenttitel rechts -----------------
% \sectionmark wird von jedem \section-Aufruf mit dem Titel als Argument
% aufgerufen; die Umdefinition setzt Nummer und Titel zusammen, sodass die
% Kopfzeile links "4 Beurteilung des Verlaufs des Aktienkurses" zeigt.
% \section* (Quellenverzeichnis usw.) loest \sectionmark nicht aus - die
% Kopfzeile behaelt dort den letzten nummerierten Abschnitt; deshalb wird
% sie vor dem Anhang manuell geleert (siehe hausarbeit.tex). \titlepage
% setzt bereits \thispagestyle{empty}, das Deckblatt bleibt also
% kopfzeilenfrei. Der rechte Eintrag ist statisch und benennt das Dokument.
\usepackage[automark]{scrlayer-scrpage}
\pagestyle{scrheadings}
\renewcommand{\sectionmark}[1]{\markright{\thesection\quad #1}}
\clearpairofpagestyles
% \small: Abschnittstitel plus Dokumenttitel muessen in eine Zeile passen.
\setkomafont{pagehead}{\normalfont\small}
\ihead{\rightmark}
% Der rechte Eintrag nennt den Teil, in dem der Leser steht. Bei einem
% Dokument mit zwei Titeln ist das keine Verzierung: Zwei gleich gebaute
% Abschnittsfolgen sind sonst auf der einzelnen Seite nicht zu unterscheiden.
\newcommand{\teiltitel}{Anlageanalyse}
\ohead{\teiltitel}
\KOMAoptions{headsepline=true}
\cfoot*{\pagemark}

% --- Inhaltsverzeichnis: nur bis subsection (x.x) --------------------
% scrartcl zaehlt section=1, subsection=2, subsubsection=3. tocdepth=2
% laesst subsubsections (z. B. 5.1.2) aus dem Verzeichnis, ohne ihre
% Numerierung im Fliesstext zu aendern - secnumdepth bleibt unangetastet,
% da die Nummern der Aufgabenstellung (5.1.1-5.1.3 usw.) im Text stehen
% bleiben muessen.
\setcounter{tocdepth}{2}

\pgfplotsset{compat=1.18}
% dateplot stellt "date coordinates in=x" bereit. Ohne die Bibliothek
% versucht pgfplots, "2020-12-31" als Gleitkommazahl zu lesen und bricht ab.
\usepgfplotslibrary{dateplot}
% pgfplots formatiert seine Achsen- und Datenbeschriftungen unabhaengig von
% siunitx. Ohne diese Zeile stuenden im Diagramm Punkte, im Text Kommata.
\pgfplotsset{/pgf/number format/use comma, /pgf/number format/1000 sep={}}

% Nach unten weisendes Dreieck als eigene Markenform. Der naheliegende Weg,
% "mark=triangle*" mit "mark options={rotate=180}" zu drehen, funktioniert
% innerhalb von scatter/classes NICHT: pgfplots setzt mark options je Klasse
% neu und verwirft die Drehung stillschweigend - die Tiefstkurse wurden
% dadurch mit einem nach OBEN weisenden Dreieck gezeichnet. Eine eigene
% Markenform kann nicht ueberschrieben werden.
\pgfdeclareplotmark{dreieckab}{%
  \pgfpathmoveto{\pgfqpoint{-0.6\pgfplotmarksize}{0.5\pgfplotmarksize}}%
  \pgfpathlineto{\pgfqpoint{0.6\pgfplotmarksize}{0.5\pgfplotmarksize}}%
  \pgfpathlineto{\pgfqpoint{0pt}{-0.55\pgfplotmarksize}}%
  \pgfpathclose%
  \pgfusepathqfillstroke%
}

% Farben der drei Punktarten in fig:kursverlauf. Gruen fuer das Jahreshoch,
% Rot fuer das Jahrestief, Blau fuer den Schlusskurs. Die Farbe kommt zur
% Form hinzu und ersetzt sie nicht: Im Schwarzweissdruck bleiben die drei
% Arten ueber Kreis, Dreieck-auf und Dreieck-ab unterscheidbar.
\definecolor{kurshoch}{RGB}{0,120,60}
\definecolor{kurstief}{RGB}{175,30,35}
\definecolor{kursschluss}{RGB}{25,55,120}

% Achsenvorlage der Kursabbildung. Sie steht hier und nicht in der
% Abbildung selbst, weil scripts/check_literals.py in sections/*.tex jede
% Zahl anschlaegt und die Tickweite von 365 Tagen ein Satzparameter ist,
% keine Finanzkennzahl. Die Regel bleibt damit unangetastet.
\pgfplotsset{
  kursachse/.style={
    width=0.95\linewidth, height=6.2cm,
    date coordinates in=x,
    xticklabel style={rotate=45, anchor=north east, font=\footnotesize},
    xticklabel=\year,
    xtick={2017-01-01,2018-01-01,2019-01-01,2020-01-01,2021-01-01,
           2022-01-01,2023-01-01,2024-01-01,2025-01-01,2026-01-01},
    ylabel={Kurs in USD}, ylabel style={font=\footnotesize},
    yticklabel style={font=\footnotesize},
    ymin=0,
    grid=both, grid style={gray!22, very thin},
    axis lines=left, enlarge x limits=0.03},
  % Zwei Titel, zwei Groessenordnungen: 166 USD gegen 12 USD. Eine
  % gemeinsame Achse waere unlesbar, und eine gemeinsame Abbildung waere
  % ein Vergleich, den dieser Bericht ausdruecklich nicht anstellt.
  kursachseazn/.style={kursachse, ymax=230},
  kursachserdy/.style={kursachse, ymax=18},
}

% Zuordnung ueber die Spalte "art" der Datendatei (S/H/T), nicht ueber
% getrennte Datensaetze - so bleibt die Reihe eine einzige, nach Datum
% geordnete Linie.
\pgfplotsset{
  kurspunkte/.style={
    scatter, only marks, scatter src=explicit symbolic,
    scatter/classes={
      S={mark=*,         mark size=1.9pt, kursschluss},
      H={mark=triangle*, mark size=2.6pt, kurshoch},
      T={mark=dreieckab, mark size=2.6pt, kurstief}}},
}

% --- Zahlenformat: deutsches Dezimalkomma ---------------------------
\sisetup{
  output-decimal-marker = {,},
  group-separator       = {.},
  group-minimum-digits  = 4,
  % Ohne diese Zeile gruppiert siunitx auch die NACHKOMMASTELLEN: aus dem
  % Rohwert 43.4679 wurde im Satz "43,467.9" statt "43,5". Die Rohwerte in
  % data/kennzahlen.tex tragen volle Genauigkeit, deshalb trat der Fehler
  % bei jeder ungerundeten Prozentangabe auf.
  group-digits          = integer,
  detect-weight         = true,
  detect-family         = true
}

% --- Betragsmakros --------------------------------------------------
% Die Rohwerte stehen in data/kennzahlen.tex im englischen Format und in
% voller Genauigkeit. Die Formatierung erfolgt ausschliesslich hier.
%   \MioUSD       - Fliesstext, auf eine Nachkommastelle gerundet
%   \MioUSDexakt  - Rechenwege und Tabellen, volle veroeffentlichte Genauigkeit
% Der Konzernabschluss weist in Mio. USD ohne Nachkommastelle aus; die
% Rundung liegt trotzdem in der Ausgabe und nicht im gespeicherten Wert,
% damit abgeleitete Groessen (Margen, Veraenderungsraten) aus den exakten
% Werten hervorgehen und nicht aus vorgerundeten.

% --- Zitierweise ------------------------------------------------------
% \Q{Schluessel}{Seite}: Fussnote mit Kurztitel (Sprung ins Quellenverzeichnis)
% und GEDRUCKTER Seite. Seite "Deckblatt" fuer SEC-Deckblaetter ohne Seitenzahl.
% \QW{Titel}{Schluessel}{Abrufdatum}: Webquelle; die Adresse steht nur im
% Quellenverzeichnis, damit jede URL genau einmal im Dokument vorkommt.
\definecolor{quellenlink}{RGB}{20,60,120}
\newcommand{\quellensprung}[2]{\hyperref[#1]{\textcolor{quellenlink}{#2}}}
\newcommand{\kurzOxyZK}{Occidental Petroleum, Form 10-K 2025}
\newcommand{\kurzOxyPS}{Occidental Petroleum, Proxy Statement 2026 (DEF 14A)}
\newcommand{\kurzTteZF}{TotalEnergies, Form 20-F 2025}
\newcommand{\kurzTteURD}{TotalEnergies, Universal Registration Document 2025}
\newcommand{\kurzTnkZF}{Teekay Tankers, Form 20-F 2025}
\newcommand{\kurzFroZF}{Frontline, Form 20-F 2025}
\newcommand{\kurzGeaGB}{GEA Group, Annual Report 2025}
\newcommand{\kurzBesiGB}{BE Semiconductor Industries, Annual Report 2025}
\newcommand{\Q}[2]{\footnote{\quellensprung{quelle:#1}{\csname kurz#1\endcsname}, S.~#2.}}
\newcommand{\QW}[3]{\footnote{\quellensprung{#2}{#1} (abgerufen am #3).}}

% --- Einheiten: Makros aus data/kennzahlen.tex sind bereits deutsch formatiert
\newcommand{\Mio}[2]{#1\,Mio.\,#2}
\newcommand{\je}[2]{#1\,#2}
\newcommand{\Pz}[1]{#1\,\%}

% --- Drei Farben, drei Verantwortlichkeiten (wie Pharma_AZN_RDY) ------
% Schwarz: belegte Fakten. Blau: Einschaetzung des Verfassers. Rot: Votum.
\definecolor{vdrot}{RGB}{155,25,30}
\newcommand{\vd}[1]{{\color{vdrot}#1}}
\definecolor{bkblau}{RGB}{20,60,150}
\newcommand{\bk}[1]{{\color{bkblau}#1}}

% --- Generierte Tabellenkoerper ohne angehaengtes \relax --------------
\makeatletter
\newcommand{\tabellenkoerper}[1]{\@@input #1.tex }
\makeatother

\newcommand{\annahme}[1]{%
  \par\medskip\noindent\fbox{\parbox{0.97\linewidth}{\small\textbf{Annahme:} #1}}\par\medskip}

% --- Drei Paare als Teile; Abschnittszaehlung je Teil neu -------------
\let\origpart\part
\renewcommand{\part}[1]{\clearpage\setcounter{section}{0}\origpart{#1}}

% --- Die fuenf Tabellen je Titel (Koerper aus scripts/rechnung.py) ----
% #1 Kuerzel (oxy, tte, ...), #2 Waehrung, #3 Label-Suffix
\newcommand{\reihentabelle}[3]{%
\begin{table}[htbp]\centering\small
\caption{Mehrjahresreihe (Mio.~#2; Ergebnis je Aktie in #2). Freier Mittelzufluss =
operativer Mittelzufluss + Investitionen + Leasingtilgung (Methode im Rahmen,
S.~\pageref{sec:methode}). Beleg: Bericht und gedruckte Seite der Kapitalflussrechnung;
\enquote{XBRL} = SEC-companyfacts, Sekund\"arquelle.}\label{tab:reihe-#3}
\begin{tabular}{@{}lrrrrrrl@{}}\toprule
Jahr & Operativ & Investitionen & Leasing & Freier Zufluss & Dividenden & EPS verw. & Beleg\\\midrule
\tabellenkoerper{data/#1_reihe}
\bottomrule\end{tabular}\end{table}}
\newcommand{\schwellentabelle}[3]{%
\begin{table}[htbp]\centering\small
\caption{Schwellenkurs je Aktie (#2): Kurs, zu dem der interne Zinsfuss die H\"urde von
\Pz{\Huerde} erreicht; letzte Spalte bei \Pz{\HuerdeLang}. Zufluss konstant, Endwert =
Buchwert des Stammeigenkapitals 2025. n.\,b.\ = nicht bewertbar (negativer Zufluss).}\label{tab:schwelle-#3}
\begin{tabular}{@{}lrrrrr@{}}\toprule
Szenario & Zufluss (Mio.) & 5 Jahre & 10 Jahre & 15 Jahre & 10 J., \Pz{\HuerdeLang}\\\midrule
\tabellenkoerper{data/#1_schwelle}
\bottomrule\end{tabular}\end{table}}
\newcommand{\umkehrtabellen}[1]{%
\begin{table}[htbp]\centering\small
\caption{Umkehrungen zum heutigen Kurs: erforderliches Wachstum des freien Zuflusses
p.\,a.\ (oben) und erforderlicher Endwert als Vielfaches des Buchwerts (unten), je bei
\Pz{\Huerde}.}\label{tab:umkehr-#1}
\begin{tabular}{@{}lrrr@{}}\toprule
Szenario & 5 Jahre & 10 Jahre & 15 Jahre\\\midrule
\tabellenkoerper{data/#1_wachstum}
\midrule
\tabellenkoerper{data/#1_endwert}
\bottomrule\end{tabular}\end{table}}
\newcommand{\kgvtabelle}[2]{%
\begin{table}[htbp]\centering\small
\caption{Kurs-Gewinn-Verh\"altnis zum Jahresende: Schlusskurs (Yahoo Finance) durch
verw\"assertes Ergebnis je Aktie (Tabelle~\ref{tab:reihe-#1}); nur Jahre mit positivem
Ergebnis.}\label{tab:kgv-#1}
\begin{tabular}{@{}lrrr@{}}\toprule
Jahr & Kurs Jahresende (#2) & EPS verw. (#2) & KGV\\\midrule
\tabellenkoerper{data/#1_kgv}
\bottomrule\end{tabular}\end{table}}
\newcommand{\paartabelle}[2]{%
\begin{table}[htbp]\centering\small
\caption{Paarvergleich: KGV heute, Median und Rang in der eigenen Reihe; freier Zufluss
2025 im Verh\"altnis zum B\"orsenwert; B\"orsenwert zu Buchwert.}\label{tab:paar-#1}
\begin{tabular}{@{}lrrrrr@{}}\toprule
Titel & KGV heute & KGV-Median & Rang (\%) & Zufluss/B\"orsenwert & KBV\\\midrule
\tabellenkoerper{data/#2}
\bottomrule\end{tabular}\end{table}}
